\section*{Declarations}

\paragraph{Funding.}
No grant or external funding supported this work. The participating administration provided
in-kind institutional support: access to its data lake and to staff who explained the meaning of
organisational units and action types.

\paragraph{Ethics and data-processing authorisation.}
No data were collected from human participants. Staff of the participating administration were
consulted, in the course of the authors' ordinary working contact with them, on what particular
organisational units and action types signify in the source system; no interview was recorded,
transcribed or analysed, no personal data were collected from them, and no statement by any
member of staff is reported as a finding. On that basis the study was not submitted to an ethics
committee. Readers who take a broader view of what constitutes research involving human
participants should treat that consultation as the point at issue, and it is disclosed here for
that reason.

The study was conducted under a formal data-processing authorisation issued by the head of the
participating administration acting as data controller under the Brazilian General Data
Protection Law, in force for the period covering this work. The instrument authorises processing of
administrative case records for academic purposes, requires pseudonymisation or anonymisation
before the clustering and vectorisation stages, requires that processing occur entirely within
the administration's own computing environment using locally executed open-weight models with
no transfer to external services, and restricts academic disclosure to cluster names, generic or
synthetic process examples, aggregate metrics, and source code without credentials. All results
reported here fall within that scope.

\paragraph{Positionality and competing interests.}
The authors have no financial competing interests. Two non-financial ones are declared. The
research team has an institutional relationship with the administration studied, which bears
both on the access that made this work possible and on the potential for bias in framing and
interpretation. And the administration held a right of prior clearance over publication of the
aggregate results reported here, which is a competing interest rather than a form of support:
findings unfavourable to the administration are structurally less likely to survive to
publication than favourable ones, and a reader should weigh the negative results reported
throughout this paper against the fact that they were published with its consent. That
clearance was sought and granted for this manuscript; the paper reports results the
administration was free to withhold, including the ones least favourable to the case for
adopting the method. A positionality
statement setting out both has been provided to the editors separately; it is withheld here
because it would identify the authors during double-blind review, and it will appear in full at
camera-ready.

\paragraph{Legal basis for the processing.}
The authorisation described above is the instrument under which this study proceeded, but an
authorisation issued by a controller is not itself a basis for processing under the general data
protection statute. The instrument does allocate responsibility for the stage that matters: it
provides that the exploratory graph may contain non-anonymised data, that this stage forms part
of the processing carried out under the controller's responsibility and confined to its secure
environment, and that it does not constitute autonomous academic use. What it does not do is
identify which statutory hypothesis the controller relied on, and neither do we. We record that
gap rather than paper over it: an administration seeking to reproduce this work would need to
establish that basis before beginning, and Section~\ref{sec:governance} states why we regard it
as the first question such an administration should answer.

\paragraph{Data availability.}
Under the governing instrument, the methodology and source code are the researcher's property
and the derived graph database is the administration's. Source code and derived aggregate
artefacts --- cluster-level metrics, the generated result tables, and the evaluation pipeline
--- can therefore be released. The underlying event log, case documents and graph database
cannot: they contain personal data whose transfer to third parties the instrument forbids, and
they are not the researcher's to release. Every quantitative claim in this paper is regenerated
from stored artefacts by a single pipeline stage, which fails if the manuscript and the
artefacts disagree.

\paragraph{Use of AI tools.}
Large language models were used in two distinct roles, both disclosed. As components of the
method: a locally executed open-weight model produced the topic and summary of each opening
document and the controlled keyword vocabulary, and a local embedding model produced document
vectors; the governing instrument requires that such models run locally, and no case data was
sent to any external service. As adversarial review assistance: AI assistance
was used to review the manuscript text and the analysis and test code, particularly to identify
logical gaps --- claims unsupported by the evidence cited, figures disagreeing between prose and
artefact, and missing steps in an argument. That use is distinct from the previous one: it
produced objections to the text rather than text. Several claims made in earlier drafts of this
paper were withdrawn or restated as a result. Every objection was assessed individually, and those accepted were accepted on verification against
the artefacts rather than by deference to the tool. The authors verified every reported figure
against the generated artefacts and are responsible for all claims.

\paragraph{Author contributions.}
Withheld for double-blind review; to be supplied in CRediT taxonomy form at camera-ready.
